%% copyright \dash{} strona redakcyjna, na odwrocie strony tytułowej.
%%
%% Stoi na stronie ii, którą \cleardoublepage i tak zostawiała pustą, więc
%% książka zyskała stronę redakcyjną bez kosztu w stronach. Wszystko, co na
%% niej stoi, było wcześniej złożone szarością u dołu strony tytułowej, czyli
%% tam, gdzie strona redakcyjna trafia, kiedy książka jej nie ma.
%%
%% \bookedition i \bookyear są przypięte w preamble.tex, a \pinnedon jest
%% łańcuchem widocznym dla czytelnika i mieszka w lang/<lang>.tex. Żadnego
%% z tych trzech nie wypisuj tu ręcznie: numer wydania i rok praw autorskich
%% to po jednym fakcie, a data przeliczenia to fakt od nich różny.
%%
%% Brak numeru ISBN jest nazwany, a nie zostawiony do wywnioskowania. Nie da
%% się go wyprowadzić, policzyć ani napisać -- jest nadawany -- więc jest tą
%% jedną rzeczą, którą ta książka jest winna czytelnikowi, a której żadne
%% przejście po tym repozytorium nie zamknie. Dodatek F wymienia go z resztą.

\thispagestyle{empty}
\vspace*{\fill}

{\small\color{mgrey}\raggedright\setlength{\parindent}{0pt}%
%  Every paragraph here is flush left. \parindent is zeroed rather
%  than left to \raggedright, which does zero it -- but this page has
%  six paragraphs and only the first would have carried a \noindent.
\emph{Matematyka od zera dla inżyniera AI}. Wydanie \bookedition.

\medskip
Copyright \textcopyright{} \bookyear{} Konrad Cinkusz. Wszelkie prawa
zastrzeżone poza tym, co udostępniają licencje poniżej.

\medskip
Tekst jest na licencji CC BY-NC-SA 4.0. Kod \dash{} każdy skrypt w katalogu
\texttt{code/} i każdy wydruk, który książka z niego bierze \dash{} jest na
licencji MIT.

\medskip
Nakładem własnym autora. Nie ma wydawnictwa, nie ma dystrybutora, nie ma
numeru ISBN ani odnotowanego miejsca wydania, więc tę książkę wolno czytać i
kopiować na powyższych licencjach, a nie można jej jeszcze sprzedawać w
księgarni ani skatalogować w bibliotece. Ten brak jest odnotowany w
Dodatku~\ref{app:F} razem z resztą tego, co ta książka jest winna
czytelnikowi.

\medskip
Wydana w dwóch językach i dwóch formatach: po angielsku i po polsku, w
$17 \times 24$~cm oraz na A4. Dostępna także po angielsku jako
\emph{Mathematics from Zero for the AI Engineer}.

\medskip
Liczby przeliczone \pinnedon. Każda liczba w tej książce pochodzi ze skryptu
w katalogu \texttt{code/} i jest sprawdzana przy każdym budowaniu. To inna
data niż rok praw autorskich powyżej i znaczy co innego.

\medskip
Metoda nauczania programowanego pochodzi od K.~A. Strouda. Ta książka nie
jest z nim ani z jego podręcznikiem powiązana.
\par}

\cleardoublepage
